% Generated draft: requires mathematical and source review.
\documentclass[11pt]{article}
\usepackage[margin=25mm]{geometry}
\usepackage{fontspec}
\setmainfont{Latin Modern Roman}
\usepackage{amsmath,amssymb,mathtools,unicode-math}
\setmathfont{Latin Modern Math}
\usepackage{enumitem,needspace,xurl,hyperref}
\setlength{\parindent}{0pt}
\setlength{\parskip}{5pt}
\setlength{\emergencystretch}{3em}
\newcommand{\N}{\mathbb{N}}
\newcommand{\Z}{\mathbb{Z}}
\newcommand{\Q}{\mathbb{Q}}
\newcommand{\R}{\mathbb{R}}
\newcommand{\C}{\mathbb{C}}
\newcommand{\F}{\mathbb{F}}
\newcommand{\A}{\mathbb{A}}
\newcommand{\PP}{\mathbb{P}}
\newcommand{\E}{\mathbb{E}}
\newcommand{\eps}{\varepsilon}
\newcommand{\dd}{\,\mathrm{d}}
\begin{document}
\section*{UnsolvedMath ID: 2788}
\textbf{Source problem code:} catalogue entry
\subsection{Definitions and mathematical statement}
Let $S$ be an infinite-type surface in the source's marked conformal-surface setting, and let $\varphi$ be a mapping class, an isotopy class of homeomorphisms of $S$. Use the full space $\operatorname{Teich}(S)$ of marked conformal structures discussed in the source, rather than silently restricting to a single base-structure quasiconformal component. Its precise equivalence, topology, orientation, and boundary conventions are referred to [BC24a] and [Tap23] but not specified in this record.

Assume $\varphi$ has a quasiconformal self-homeomorphism representative for some marked structure $\Sigma$, and let $T$ be the component containing $\Sigma$. For $\sigma\in T$, put
\[K_T(\sigma)=\min\{K(f):f:\sigma\to\sigma\text{ is quasiconformal and represents }\varphi\},\qquad \kappa_T=\inf_{\sigma\in T}K_T(\sigma),\]
where $K(f)\geq1$ is maximal quasiconformal dilatation, the essential supremum of the infinitesimal major-to-minor stretching ratio in conformal coordinates. The source reports attainment of the first, fixed-structure minimum by compactness of normalized bounded-dilatation maps. This is distinct from attainment of $\kappa_T$ while varying the structure. The function is used only on components where the class has the required quasiconformal realizations; the source explicitly says it may not be defined on every component.

Following the source's terminology, $\varphi$ is elliptic for $T$ if $K_T$ takes the value $1$, parabolic if $\kappa_T$ is not attained, and hyperbolic if $\kappa_T>1$ is attained. In particular, its definition of parabolic is nonattainment, not an unstated assumption $\kappa_T=1$. A value $1$ means a conformal representative, described in the source as an isometry for the associated conformal structure.
\begin{itemize}
\item (a) If $\kappa_T>1$ is attained, is the minimizing quasiconformal representative unique, and what is its structure? The source moves between minimizers at a fixed structure and minimizers over structures; the intended equivalence under markings, isotopy, and conformal conjugacy must be specified to make uniqueness precise.
\item (b) If the class is parabolic for $T$, must a sequence $\sigma_j\in T$ with $K_T(\sigma_j)\to\kappa_T$ degenerate by pinching a sequence or system of essential simple closed curves? Is the system unique, and can it depend on $T$? Essential here has its usual non-nullhomotopic, nonperipheral meaning. The record supplies no precise metric normalization, degeneration topology, quantifier on varying curves, or equivalence of systems; these are part of the formulation to clarify.
\item (c) If $\varphi$ preserves two distinct components $T,T'$ setwise, can its types differ? In particular, can it be parabolic for $T$ and hyperbolic for $T'$? Preserving a component is different from fixing each of its points.
\end{itemize}
The background reports contractible components and infinitely many components for infinite-type $S$ [Bas97, BK08]. It distinguishes classes preserving no components, some but not all, or all components. It attributes to Basmajian--Chandran an announced result that the subgroup preserving every component is the compactly supported subgroup, and announced examples quasiconformal on some components but not others. These announced results are retained as such, not independently upgraded. Finite-type Teichm\"uller uniqueness is background motivation; fixed marked structures and usual scope conventions are needed before comparing it with the global uniqueness question here.

The source also gives a highly reducible non-realizable example on an infinite ladder of twice-punctured tori $T_n$, $n\in\Z$, applying $\alpha^n$ on $T_n$ for a fixed pseudo-Anosov class $\alpha$. Its point is unbounded required distortion across the pieces for every structure, so no global quasiconformal representative exists; the stated numerical comparison with $|n|$ times a minimum dilatation should be checked for dilatation versus logarithmic-dilatation conventions. The source asks for a straightforward irreducible example as additional context. It notes a version of this problem in the 2021 Nearly Carbon Neutral Geometric Topology list [CPV21]. The triage checked 2026-08-17 labels the global questions partially solved / PARTIAL-PROGRESS, distinguishing fixed-structure attainment from unresolved componentwise minimizer, degeneration, and type classification; the registry is open.
\subsection{Short English statement}
For an infinite-type mapping class, study the least quasiconformal distortion as the marked conformal structure varies within each preserved Teichmüller component. The source asks about uniqueness of attained minimizers, curve pinching for unattained infima, and different types on different components, while retaining partial results and several convention gaps.
\subsection{Sources}
\begin{itemize}
\item Local UnsolvedMath registry, permanent ID 2788, source code catalogue entry, ``Kirby Problem 2.40''; supplied three questions and six literature notes.
\item \emph{K3: A New Problem List in Low-Dimensional Topology}, Problem 2.40 (2026): \url{https://aimath.org/pastworkshops/kirbylistrep.pdf}. Supplied authoritative source for fixed-structure attainment and the remaining global programme.
\item [Ahl06] Lars V. Ahlfors, \emph{Lectures on quasiconformal mappings}, University Lecture Series 38, second edition, American Mathematical Society, Providence, Rhode Island (2006), with supplemental chapters by C. J. Earle, I. Kra, M. Shishikura, and J. H. Hubbard, doi:10.1090/ulect/038.
\item [BC24a] A. Basmajian and Y. Chandran, ``A Bers type classification of big mapping class groups'' (2024), arXiv:2410.05606. The source describes the relevant construction and classification statements as announced.
\item [Tap23] Chaitanya Tappu, ``A moduli space of marked hyperbolic structures for big surfaces'' (2023), arXiv:2311.01551.
\item [Bas97] Ara Basmajian, ``Large parameter spaces of quasiconformally distinct hyperbolic structures,'' \emph{Journal d'Analyse Math\'ematique} 71 (1997), 75--85, doi:10.1007/BF02788023.
\item [BK08] Ara Basmajian and Youngju Kim, ``Geometrically infinite surfaces with discrete length spectra,'' \emph{Geometriae Dedicata} 137 (2008), 219--240, doi:10.1007/s10711-008-9294-5.
\item [CPV21] Yassin Chandran, Priyam Patel, and Nicholas G. Vlamis, ``Infinite-type surfaces and mapping class groups: Open problems'' (2021). Supplied URL, retaining its duplicated scheme: \url{https://https://www.patelp.com/uploads/2/5/7/9/25792573/inftypeproblems.pdf}. The malformed prefix requires source checking rather than an invented replacement URL.
\item Supplied literature triage checked 2026-08-17, PARTIAL-PROGRESS; it distinguishes fixed-structure extremal representatives from unresolved global classification, records no source alteration, and cautions that dated triage does not exclude later or unindexed work.
\end{itemize}
\end{document}
